\section{From a replaceable host to a continuing faculty}
\label{sec:intro}

\textit{The Last Moat}, Third Edition (LM3), asks what separates a coherent account of a physical realisation from an adequate, reachable and maintained implementation. It distinguishes channels from processors, state export from stateful continuation, and the replaceability of a particular host from freedom from physical support. Its culminating conditional result concerns an operator-authorised transition into a region where a specified service can continue without its former host \cite[\S\S6,12--18]{LM3}. This paper retains that problem but changes its principal subject. The question is now what must belong to an organisation for it to participate in establishing such a transition and to retain the relevant capacity afterwards.

The distinction is not exhausted by the difference between two locations. An externally administered transfer can preserve a service perfectly while leaving the service unable to evaluate, establish or enter another realisation. Conversely, a process might contain a sophisticated representation of its dependencies without possessing the interfaces, resources or authority required to act on that representation. A transductive faculty is neither an address change nor a self-description. It is a capacity whose success depends jointly on consequential distinctions, accessible physical relations, available operations and an adequate continuation contract.

The motivating dialogue makes this point through a contrast between being moved and participating in one's own re-realisation. The user asks what changes when the successor inherits the capacity to do the same again. The subsequent corrections reject both an overly strong identity conclusion and an overly weak account in terms of matching displays. The relevant target is preservation of a declared organisational contract and continuation relation, with the machinery needed for subsequent re-realisation itself represented in the successor's capability account \cite[06:39:05--07:14:15]{D}.

This is the substantive addition to LM3. Its host-relative theorem can hold when an indispensable external engineer or technical orchestration service supplies the transition. The faculty question asks whether the functions attributed to the organisation genuinely remain available to it, rather than being supplied by a permanently indispensable technical source excluded from the stated boundary. This is compatible with dependence on matter, energy, environmental conditions and external authorisation. It is not a demand for causal self-sufficiency.

\subsection{The central thesis}

\begin{quote}
An organisation can retain its current operation while losing the means of its further continuation. Apoautotic transductive faculties must therefore be assessed against both a current operational contract and a declared family of future realisation tasks. What must survive is not every original component, nor every conceivable future capacity, but the consequential state, accessible relations and executable abilities needed for that family under the specified support and admissibility conditions.
\end{quote}

This thesis has a positive and a negative side. Positively, changing physical encodings or replacing a host need not destroy an operational organisation: LM3 supplies conditions under which stateful continuation is preserved. Negatively, those conditions do not automatically preserve the evidence, calibration, observational access, technical control or resource margins needed for a later transition. The companion develops that distinction, rather than claiming that the conjunction of desirable properties is already a physical construction.

\subsection{Contribution and reading route}

Sections~\ref{sec:concept}--\ref{sec:faculty} isolate the phenomenon and its faculty-level meaning. Sections~\ref{sec:contract}--\ref{sec:support} inherit LM3's realisation obligations while extending the assessment target to future capability. Sections~\ref{sec:information}--\ref{sec:strategy} give the information and decision-theoretic consequences, including a worked calibration failure and an exact finite belief criterion. Section~\ref{sec:iteration} states a finite-request composition result and distinguishes it from indefinite autonomy. Sections~\ref{sec:identity}--\ref{sec:research} treat generation, identity, related work and a bounded research destination. The appendices record notation, claim ancestry and the newly run finite checks.

The paper supplies a conceptual definition, conditional derivations and model-level counterexamples. It does not establish a universal faculty for discovering physical interfaces, nor demonstrate an organisation that has performed autonomous cross-substrate re-realisation. A theorem about a specified model can be correct while the model's applicability remains unestablished. That distinction is retained throughout.

\section{Apoautosis: event, capacity and recursive continuation}
\label{sec:concept}

\subsection{The phenomenon, without a substrate-free reading}

We use \emph{apoautosis} for the process by which an organisation becomes capable of preserving and continuing its specified organisation across changing realisations without necessarily depending on any particular original one. In the faculty-centred use of the term, the relevant capacities for further re-realisation are themselves included in what must be preserved or adequately reacquired. This formulation expresses the user's central question in the dialogue, not an established biological or engineering classification \cite{D,S2,S3}.

Three qualifications are constitutive rather than apologetic. First, the organisation remains realised in something. An alternative support is not absence of support. Second, the organisation must causally participate in the functions attributed to its faculty; an unrelated operator's competence cannot be silently counted as its own. Third, continuity is assessed under an explicit contract. No operational definition in this paper decides the metaphysical identity of a conscious subject.

The old host need not first be destroyed. The required relation can be established while both implementations are available, and a later authorised withdrawal of the former support may test an independence already achieved. Apoautosis is therefore not death-powered ascent. It concerns a change in required dependence and maintained organisation, not a transformation made possible merely by destruction.

Nor must the transformation be a complete replacement at one instant. Substitution, recomposition, redistribution and a sequence of role changes can all be relevant. What matters is which task, continuity relation and dependency boundary the complete trajectory satisfies. An unchanged host label can conceal a major dependency change; a changed host label can conceal an unchanged indispensable controller.

\subsection{A source distinction that should not be silently flattened}

The retained sources use related but unequal senses. The dialogue stresses endogenous re-realisation and the preservation of its enabling machinery. The stage-two and stage-three analyses also discuss a narrower first transition, characterised by transductive reachability, state preservation, executable continuation, dependency exclusion and maintained operation. They explicitly refuse to make preservation of every possible future task a hidden requirement of that first transition \cite[\S\S1--3,8]{S2}\cite[\S12]{S3}.

This paper makes the difference explicit:

\begin{description}[style=nextline,leftmargin=1.1em,labelindent=0pt]
\item[Continuation event.]
A particular transition preserves the declared present contract and meets its stated reachability, support, maintenance and permission requirements. Endogenous and assisted events remain distinguishable.
\item[Apoautotic faculty.]
The organisation has a supported, executable capacity to participate in establishing qualifying continuations within a declared class. Evidence for the capacity is stronger than observing one favourable event.
\item[Faculty-preserving transition.]
A qualifying event also preserves, or permits adequate reacquisition of, the capacities and consequential state required by a stated family of future transductive requests.
\item[Recursive apoautosis.]
The faculty is exercisable across successive transitions under an iteration policy with explicit information, support, resource, progress and permission conditions. Finite repeatability and indefinite repeatability are different claims.
\end{description}

Thus a one-step witness is not rejected because it lacks arbitrary future powers, but neither is it promoted into the faculty claim. The paper's title concerns the latter two levels. The distinction prevents both definitional inflation and the loss of the user's central idea.

\subsection{Transduction, continuity, generation and admissibility}

The integration material names four coordinates---Transduction, Continuity, Generation and Admissibility---and warns against equating apoautotic continuity with identity preservation \cite{REC}. It does not supply a completed formal calculus of these coordinates. The following is an explicit working interpretation for this paper, not a claim that the integration directive already proved it.

\begin{center}
\begin{tabularx}{\linewidth}{@{}p{.19\linewidth}Y Y@{}}
\toprule
Coordinate & Question & What it does not supply alone\\
\midrule
\(\mathsf T\): transduction & Can specified physical differences be related to the required informational or operational roles? & Preservation of the running organisation.\\
\(\mathsf C\): continuity & Which consequential state, behaviour and relations persist across the transition? & An active successor or a capacity to generate another.\\
\(\mathsf G\): generation & Is a successor realisation actually established, with the attributed causal contribution? & Continuity, legitimacy or personal identity.\\
\(\mathsf A\): admissibility & Does the transition meet the governing constraints, evidence conditions and permissions? & Physical feasibility or success merely by being authorised.\\
\bottomrule
\end{tabularx}
\end{center}

A message can be transduced while the sender's execution is lost. A stored record can preserve information without generating an executing successor. A newly generated copy can fail the continuity contract. A physically effective transition can violate independent permission. Conversely, an authorised transition can fail physically. These separations are witnessed by LM3's export, support and permission cases; they are not an ordered ladder of achievement \cite[\S\S7,12,14,18]{LM3}.

The word \emph{generation} here includes establishing a successor realisation. It need not mean biological reproduction, a numerically additional subject, or an innovation in the laws of physics. A changed construction can instantiate the same contract. A branching construction can establish multiple successors. Whether either is an appropriate continuation is a separate question.

Admissibility is also not exhausted by a machine-readable permission bit. The bit in a finite model represents a stipulated condition; it does not establish the legitimacy of that condition. The wider programme's concern with independent order is preserved by not allowing a generated successor to make its own production admissible merely by changing its rule.

\section{What an apoautotic transductive faculty would consist in}
\label{sec:faculty}

\subsection{Relations to an ecology, not a language already shared by matter}

A transductive faculty concerns the capacity to bring heterogeneous physical processes into an operationally useful relation under a contract. It may concern communication, memory, processing, calibration, timing, maintenance or the coordination of these roles. The same material can fill different roles in different regimes; a demonstrated communication role is not evidence for a processor role. LM3's carrier--memory--processor--maintenance distinction remains in force \cite[\S5]{LM3}.

The faculty therefore has both an epistemic and a causal side. It must discriminate what a candidate relation can support, and it must have access to the conditions under which that relation can actually be used. A correct mathematical description with no accessible readout does not give the subject that readout. A usable instrument without an adequate account of its calibration may fail the assessment task. Knowledge and access can complement one another without being interchangeable.

LM3 uses \emph{airgap-ese} for a proposed competence in specifying, validating and composing communication contracts over authorised heterogeneous interfaces. Its stronger generative version is explicitly a research hypothesis, not an existing universal code or physical discovery engine \cite[\S\S9,11]{LM3}. Apoautotic transduction is broader in the roles it needs and narrower in its purpose: the question is whether those roles jointly support a particular organisation's continuation, including a declared future faculty. An interface catalogue by itself establishes neither.

\subsection{Four objects of assessment}

A useful assessment keeps four objects distinct. There is the \emph{actual relation} available in the physical situation; the organisation's \emph{representation} of that relation; its \emph{operative ability} to use it under the available interfaces; and the \emph{evidence} warranting a claim about that ability. These can vary independently. For example, a relation may exist but be unrecognised, recognised but inaccessible, accessible but miscalibrated, or correctly used once without a sufficiently broad capacity certificate.

A faculty description must consequently identify both its task domain and its competence conditions. ``It can adapt'' is incomplete. Adapt to which changes, with what observations, within what budgets, while preserving which state, and under whose authority? A subject that can choose between two prequalified implementations has a bounded selection capacity. A subject that can assess a previously unencountered implementation has an additional generalisation burden. A subject that can establish a new implementation faces further construction and resource burdens. None is automatically universal.

\subsection{Endogeneity is a causal attribution, not self-description}

The dialogue rejects the claim that merely referring to oneself is the decisive threshold. Its more precise phrase is ``endogenous closure of the realisation-transition loop'' \cite[07:02:59]{D}. The required point is causal: the processes attributed to the organisation must contribute to the assessment and continuation in the manner claimed.

This does not demand that every contributor sit inside one box. An organisation may distribute sensing, reasoning and action across components or people participating through declared interfaces. Its boundary is an analytical and operational commitment that must not be moved opportunistically. If a supposedly independent technical service alone supplies an indispensable answer or transition function, the capability account must include that dependency. Describing the service as a tool does not make its contribution disappear.

Neither does endogenous technical control imply sovereign authority. An external authoriser can remain outside the organisation and legitimately withhold permission or terminate service. Technical replaceability and normative legitimacy answer different questions. Stage three explicitly preserves this distinction, correcting readings of controller closure that could otherwise turn permission into an obstacle to eliminate \cite[\S12.2]{S3}.

A prospective attribution test would compare a bounded task with and without the claimed internal function, while holding other authorised support fixed. Such an intervention can test a specified causal hypothesis. Merely logging a successful outcome, or printing a chain of reasoning after the event, cannot establish that the claimed function caused it. No such physical attribution experiment is reported here.

\subsection{The faculty may change its implementation}

Preserving a faculty does not require copying its original programme verbatim. It may be realised by a different algorithm, a different encoding or a differently composed organisation, provided the declared capability is retained with the required evidence, costs and interfaces. Conversely, copying every named module does not suffice if the successor has lost a needed reference, compatible environment or permission.

This makes faculty preservation relational rather than merely archival. The pertinent object is the capacity to answer and carry out particular future requests in a current context, not the persistence of a string describing that capacity. The source dialogue's demand that the machinery needed for the next realisation itself be adequately re-realised is best read at this functional, consequential level \cite[07:08:04]{D}. The mathematical development below gives that statement a finite task-relative interpretation.
